% Generated draft: requires mathematical and source review.
\documentclass[11pt]{article}
\usepackage[margin=25mm]{geometry}
\usepackage{fontspec}
\setmainfont{Latin Modern Roman}
\usepackage{amsmath,amssymb,mathtools,unicode-math}
\setmathfont{Latin Modern Math}
\usepackage{enumitem,needspace,xurl,hyperref}
\setlength{\parindent}{0pt}
\setlength{\parskip}{5pt}
\setlength{\emergencystretch}{3em}
\newcommand{\N}{\mathbb{N}}
\newcommand{\Z}{\mathbb{Z}}
\newcommand{\Q}{\mathbb{Q}}
\newcommand{\R}{\mathbb{R}}
\newcommand{\C}{\mathbb{C}}
\newcommand{\F}{\mathbb{F}}
\newcommand{\A}{\mathbb{A}}
\newcommand{\PP}{\mathbb{P}}
\newcommand{\E}{\mathbb{E}}
\newcommand{\eps}{\varepsilon}
\newcommand{\dd}{\,\mathrm{d}}
\begin{document}
\section*{UnsolvedMath ID: 10400092}
\textbf{Source problem code:} catalogue entry
\subsection{Definitions and mathematical statement}
Let $S_{4,\infty}(S^3;R)$ be the source's framed-link skein module over the chosen coefficient ring $R$, with its $(4,\infty)$ skein and framing parameters. Its local relation diagrams and full parameter conventions are not reproduced. For the allowed trivial unlinks $U_m$, the three questions attributed to J. Przytycki are:

(i) Find an $R$-module generating set for $S_{4,\infty}(S^3;R)$.

(ii) Determine exactly which parameter choices make the classes $[U_m]$ linearly independent, meaning every finite relation $\sum_m a_m[U_m]=0$ has all $a_m=0$.

(iii) Determine which parameter choices make the trivial knot $[U_1]$ non-torsion. Over a domain this means no nonzero coefficient annihilates $[U_1]$. For a general $R$, the intended torsion convention must be fixed; annihilator-zero and regular-coefficient torsion conventions need not be interchangeable.

The source discusses a former generalized Montesinos--Nakanishi conjecture: trivial links would generate the sphere module, and a specified set of basic $n$-tangles would generate relative modules with $2n$ boundary points in $S^3$ or $D^3$. Its proposed count is extracted as $\prod_{i=1}^{n-1}(3i+1)$, with potentially lost superscripts; the common exponential reading would be $\prod_{i=1}^{n-1}(3^i+1)$. The original typographic expression must be checked rather than silently selecting one. The source says this count is only a lower bound after the 3-move conjecture's disproof, not a sufficient generating-set bound; no such lower bound is derived merely from the logical fact of disproof here.

For the parameter subcase $b_\infty=0$, the source gives the necessary condition
\[
b_0b_1=b_2b_3
\]
for trivial-link linear independence. It is not asserted sufficient, nor is it automatically a criterion for part (iii). The extract then announces a following conjecture, cases (1)--(2), but those cases are absent. No missing parameter classification is reconstructed.
\subsection{Short English statement}
Find generators and determine the unlink-independence and unknot-nontorsion parameter loci for the $(4,\infty)$ skein module. Retain the disproof of the earlier 3-move generation proposal and the source's necessary parameter relation.
\subsection{Sources}
\begin{itemize}
\item Local UnsolvedMath registry, permanent ID 10400092, exact source code catalogue entry, Problem 4.14; source attribution J. Przytycki, no proposer/year in metadata. Registry status open, classification OPEN-TRIAGE; title truncated.
\item Ohtsuki, \emph{Problems on invariants of knots and 3-manifolds} (2002), Problem 4.14, PDF page 78, \url{https://msp.org/gtm/2002/04/gtm-2002-04-024s.pdf}. Supplied pdf-font-structure extraction accessed July 29, 2026; provenance distinguishes 231 bold numbered items. No PDF was accessed here.
\item Source reference [345] is cited for the generalized generation proposal and parameter analysis. The source credits M. Dabkowski and J. H. Przytycki, February 2002, references [99] and [342], for disproving the Montesinos--Nakanishi 3-move conjecture. Bibliography entries are absent; this historical result does not itself answer all three module questions.
\item The supplied summary repeats that disproof but supplies no precise current module classification. Provenance flags current primary-literature status and redistribution terms NEEDS\_REVIEW.
\end{itemize}
\end{document}
